\chapter{Introduction}
\label{ch:intro}

\section{Motivation}
\label{sec:motivation}

Distributed Denial-of-Service (DDoS) attacks have grown in both volume and
diversity. Recent industry reports record multi-terabit-per-second volumetric
floods and simultaneously show a rise in application-layer
and low-rate attacks that evade simple rate thresholds. In modern
cloud data-center networks (DCNs) these attacks threaten not only single
victims but also co-tenants sharing the physical substrate: a volumetric
surge against one tenant can saturate fat-tree aggregation or core links
and violate the service-level agreements (SLAs) of neighbouring tenants.

The dominant defensive posture today is \emph{reactive}: detection rules
fire only after the offending feature --- typically packet rate or SYN
ratio --- crosses a static threshold, at which point the control-plane
instantiates a mitigation Virtual Network Function (VNF). The latency of
this loop is dominated by VNF boot time. In our measurements the reactive
path required approximately \(6{,}006\) ms end-to-end; during this window
the attack proceeds unchecked.

This latency is not a property of any specific orchestrator --- it is
structural. Reducing it requires acting \emph{before} the attack crosses
the reactive trigger, which in turn requires a forecast.

\section{Problem statement}
\label{sec:problem}

We target the orchestrator of a multi-tenant DCN whose east-west and
north-south traffic is monitored via gNMI and NetFlow, whose mitigation
capacity is expressed as instantiable VNFs under an ONAP MANO control
plane, and whose tenants hold differentiated SLAs (e.g. URLLC with
guaranteed floor, eMBB best-effort, mMTC elastic). The defender must decide,
for every 5-second telemetry window, \emph{which} tier of mitigation to
apply --- ranging from do-nothing to heavy scrubbing --- so as to minimise
attack impact while preserving benign-tenant SLAs and keeping false-positive
actuation bounded.

Existing systems either (i) detect but do not orchestrate, (ii) orchestrate
reactively without forecasting, or (iii) forecast but do not integrate with
a standards-compliant MANO stack nor with a multi-tenant SLA model.

\section{Research questions and hypotheses}
\label{sec:rq}

We formalise the scope of this thesis through three research questions
and three testable hypotheses.

\begin{itemize}
  \item \textbf{RQ1 (Lead-time).} Does an ML forecaster reduce end-to-end
        mitigation latency compared with a rule-based threshold baseline
        at a controlled false-positive rate?
  \item \textbf{RQ2 (SLA).} Does a graduated five-tier policy engine
        preserve tenant SLAs under active mitigation?
  \item \textbf{RQ3 (Safety).} Is pre-positioning at tier \tier{2} safe ---
        i.e. does it not degrade benign tenants below their SLA floor, and
        can it be rolled back within bounded time?
\end{itemize}

\begin{itemize}
  \item \textbf{H1.} On gradual attacks (SYN-ramp), the ML forecaster
        triggers mitigation \(\geq 30\) s earlier than the threshold
        baseline.
  \item \textbf{H2.} Binary detection accuracy \(\geq 98\%\) and AUC
        \(\geq 0.99\) under a temporal 80/20 train/test split with no
        shuffling.
  \item \textbf{H3.} Under all eight evaluation scenarios the URLLC
        bandwidth floor of 150 Mbps is preserved at every tier.
\end{itemize}

\section{Contributions}
\label{sec:contrib}

\begin{enumerate}
  \item An \emph{AI-augmented four-tier architecture} (S1: telemetry;
        S2: features; S3: AI inference; S4: orchestration) that plugs a
        proactive forecast output into an ONAP closed loop.
  \item A \emph{multi-horizon DDoS forecaster} (Transformer+LSTM,
        \(t{+}30,60,90,120\) s) whose AUC degradation curve is consistent
        with the physical forecast-horizon trade-off.
  \item A \emph{linear-programming multi-tenant SLA allocator} that
        guarantees URLLC floors under VNF overhead up to 700 Mbps.
  \item A \emph{quantitative comparison} against a rule-based threshold
        baseline showing a 65 s lead-time advantage on gradual attacks
        and an out-of-distribution (OOD) robustness advantage (baseline
        FAIL on S4/S5; \padonap{} PASS).
\end{enumerate}

\section{Thesis outline}
The remainder of the thesis is structured as follows.
Chapter~\ref{ch:background} surveys DDoS defense, NFV/ONAP, and ML for
network security, and positions \padonap{} against related work.
Chapter~\ref{ch:architecture} presents the proposed architecture, its
threat model, and the safety invariants.
Chapter~\ref{ch:ai} details the AI components --- feature engineering,
XGBoost detection, Transformer--LSTM forecasting, validation, and
interpretability.
Chapter~\ref{ch:eval} reports the experimental evaluation against the
threshold baseline, lead-time analysis, and SLA behaviour.
Chapter~\ref{ch:conclusion} concludes and outlines future work.
